\chapter{Commenti finali}

\section{Autovalutazione e lavori futuri}

\subsection{Albini Pietro}

Avendo già molta familiarità con interpreti e parser, la struttura generale e
gli algoritmi da utilizzare mi erano già chiari in partenza, però non avevo mai
applicato queste mie conoscenze a Java. La mia difficoltà principale nel
progetto è stata infatti adattarle alla programmazione ad oggetti e agli idiomi
di Java, e sono soddisfatto del design finale che ho implementato. Verso la
fine mi sono anche ritrovato a sviluppare parti dell'interfaccia grafica per
aumentare il conto ore: non avendo mai sviluppato parti di videogiochi o
comunque di interfacce grafiche iniziare è stato un po' ostico, ma una volta
familiarizzato con i concetti non è stato difficile portare a termine il
lavoro. Essendo il membro del gruppo con più esperienza nel campo della
programmazione mi sono ritrovato ad essere il punto di riferimento per molte
scelte di design o domande tecniche, e anche se ho dato i miei suggerimenti
quando richiesti ho sempre cercato di non coprire il lavoro degli altri, ma di
portare avanti, migliorare ed enfatizzare le idee degli altri membri del
gruppo, anche nei casi in cui avrei scelto un approccio diverso se fosse tutto
dipeso da me.

\subsection{Cavalieri Giacomo}

Lavorare alla realizzazione di un progetto di queste dimensioni è stata per me
un'esperienza molto istruttiva con la quale non avevo mai avuto la possibilità
di cimentarmi prima d'ora. Lavorare in gruppo insieme ad altri studenti mi ha
permesso di apprezzare i vantaggi di avere un team affiatato a cui poter
rivolgere domande, esporre dubbi e con cui confrontarsi attivamente. Ho
ritenuto proprio il confronto scaturito dai nostri frequenti incontri la parte
più gratificante del progetto. Inizialmente l'uso del DVCS si è rivelato
particolarmente difficile per me, essendo nuovo a tale strumento; tuttavia col
passare del tempo ne ho potuto apprezzare l'utilità e ritengo sia stato di
essenziale importanza nell'ottenere un buon risultato finale. Inoltre, grazie
alle review effettuate sul mio codice e sul codice degli altri membri del
gruppo, ho avuto modo di imparare molto dai miei compagni. Durante il progetto
avevo iniziato a lavorare ad una parte della interfaccia grafica che tuttavia,
avendo già raggiunto il monte ore richiesto, ho preferito lasciare realizzare a
Pietro che aveva bisogno di aggiungere altre ore di sviluppo; per questo in
futuro mi piacerebbe provare a cimentarmi in questo campo col quale ho avuto a
che fare solo superficialmente.

\subsection{Di Domenico Nicolò}

Essendo stato responsabile di gran parte della View, è ricaduta su di me la
scelta della libreria grafica da usare. Conoscevo già libGDX, e sapevo che era
ampiamente usata per fare giochi, quindi ho deciso di dare un'occhiata a
qualche esempio e mi aveva convinto. Tuttavia, solo a progetto inoltrato mi
sono accorto che la sua architettura rende difficile una buona progettazione a
oggetti. Se avessi guardato più in dettaglio la sua Javadoc e non mi fossi
limitato al suo manuale, mi sarei accorto prima di questo suo difetto, e
probabilmente la scelta di libreria sarebbe ricaduta su qualcos'altro.

\subsection{Guerra Michele}

Il mio lavoro non riguarda gli aspetti più interessanti di programmazione
all'interno del nostro progetto, tuttavia ritengo che dopo un buon lavoro di
analisi fatto con i miei compagni io sia stato in grado di produrre una
porzione di Model facilmente estendibile, al punto da poter quasi permettere a
un giocatore con la minima conoscenza di XML di creare i propri livelli. Ho
avuto forti difficoltà in partenza di progetto per quanto riguarda l'utilizzo
del DVCS, siccome fino a quel momento la mia unica esperienza a riguardo si era
limitata a quella di laboratorio, senza mai mettere in pratica il lavoro
collaborativo. Al termine del lavoro posso dire di essere sicuramente
migliorato sotto questo aspetto, nonostante ci siano ancora molte cose che devo
imparare a riguardo.

\section{Difficoltà incontrate e commenti per i docenti}

Una delle maggiori difficoltà incontrate nel progetto è stato lavorare con la
libreria libGDX. La sua architettura tutt'altro che ottimale ci ha obbligati a
lavorare direttamente su classi e non su delle interfacce, impedendoci di
astrarre al meglio i problemi di tipo grafico che abbiamo incontrato. Sebbene
abbia alcune classi di utility molto comode (per esempio il drag and drop o
un'implementazione di un'interfaccia grafica con un sistema di layout), gran
parte della libreria è scarsamente o per niente documentata, costringendoci più
volte a dover andare a vedere il suo sorgente per capire il motivo di alcuni
comportamenti anomali, o per comprendere cosa significasse un parametro non
documentato di un metodo. Inoltre, l'interfaccia grafica ha rappresentato la
maggiore difficoltà: è stato particolarmente difficile ottenere il risultato
finale poiché, nonostante noi avessimo un'idea di partenza per un layout
grafico, trovare la combinazione di metodi giusta da chiamare per ottenerlo si
è ridotto a un mero trial and error, dato che non sempre le funzioni di layout
lo cambiavano come ci aspettavamo.

Un altro aspetto su cui ci siamo trovati in difficoltà durante lo sviluppo del
progetto è stato il calcolo delle ore spese: i membri del nostro gruppo hanno
conoscenze ed esperienze molto eterogenee, ed un'ora spesa da una persona
produce risultati differenti dall'ora spesa da altri. Per poter permettere ad
ognuno di raggiungere un monte ore accettabile abbiamo dovuto spostare parte
del carico di lavoro ai membri con più esperienza, dato che hanno completato le
loro parti obbligatorie e opzionali in molto meno tempo.
